\documentclass{article}
\title{Contributing to KDEEcoLab - structure}

\begin{document}

% turn off pagenumbering
\pagenumbering{gobble}
\maketitle

\section{Get familiar with the project}
\subsection{Read all ressources (GitLab repository, Project's handbook)}
\subsection{Run the project locally and test the scripts}
\subsection{Change the source code and perform a merge request}
\subsection{Introduction to the project and topic}

\section{Selenium}
\subsection{Evaluation of the framework for fulfilling the task}
\subsection{Learning Selenium}
\subsubsection{How does it work}
\subsubsection{Building first test scripts}

\section{Transfer scripts to the target framework Selenium for each application and test-case}
\subsection{Baseline, Idle, Standard Usage Scenario}
\subsubsection{Okular}
\subsubsection{Kate}
\subsubsection{LibreOffice}

\section{Setup Testing Lab}
\subsection{Read documentation of the software energy consumption measurement process}
\subsection{Choose Software and Hardware components and plan the setup}
\subsection{Interfaces between used components}
\subsection{Test the Lab-setup}

\section{Testing software for it's energy consumption}
\subsection{Making actual testruns for each scenario}
\subsection{Process the collected data}

\section{Evaluation}
\subsection{Evaluation of the testing scripts}
\subsection{Evaluation of the measurement results}

\newpage

\begin{center}
\section*{Rough timeline for project points:}
\begin{tabular}{ | l | r | }
 \hline
 Task & Deadline \\
 \hline\hline
 1: Get familiar with the project  & Nov \\
 \hline
 2: Selenium & Dec\\
 \hline
 3: Transfer scripts & Feb\\
 \hline
 4: Setup Testing Lab & Mar\\
 \hline
 5: Test software's energy consumption & Apr\\
 \hline
 6: Evalutaion & May\\
 \hline
\end{tabular}
\end{center}

\end{document}
